\documentclass[english,a4paper,nolineno]{lipics-v2019}
%\documentclass[english,nolineno]{socg-lipics-v2019}
%This is a template for producing LIPIcs articles. 
%See lipics-manual.pdf for further information.
%for A4 paper format use option "a4paper", for US-letter use option "letterpaper"
%for british hyphenation rules use option "UKenglish", for american hyphenation rules use option "USenglish"
% for section-numbered lemmas etc., use "numberwithinsect"
\newcommand{\R}{\mathbb{R}}

\newcommand{\nula}{\mathbf{0}}
\newcommand{\conv}{\mathop{\mathrm{conv}}}

\usepackage{microtype}%if unwanted, comment out or use option "draft"

%\graphicspath{{./graphics/}}%helpful if your graphic files are in another directory

\bibliographystyle{plainurl}% the recommnded bibstyle

\title{The Crossing Tverberg Theorem}

\titlerunning{The Crossing Tverberg Theorem}%optional, please use if title is longer than one line

\author{Radoslav Fulek}{IST Austria, Klosterneuburg,  Austria}{radoslav.fulek@gmail.com}{https://orcid.org/0000-0001-8485-1774}{The author gratefully acknowledges support from Austrian Science Fund (FWF), Project~ M2281-N35.}%mandatory, please use full name; only 1 author per \author macro; first two parameters are mandatory, other parameters can be empty.

\author{Bernd G{\"a}rtner}{Department of Computer Science, ETH Z{\"u}rich, Switzerland}{gaertner@inf.ethz.ch}{https://people.inf.ethz.ch/gaertner/}{}

\author{Andrey Kupavskii}{University of Oxford, UK; Moscow Institute of Physics and Technology, Russia}{kupavskii@ya.ru}{http://kupavskii.com}{The research was supported by the Advanced Postdoc.Mobility grant no. P300P2\_177839 of the Swiss National Science Foundation}

\author{Pavel Valtr}{Department of Applied Mathematics, Faculty of Mathematics and Physics, Charles University, Prague, Czech Republic\\ and \\ Department of Computer Science, ETH Z{\"u}rich, Switzerland}{}{http://orcid.org/0000-0002-3102-4166, https://kam.mff.cuni.cz/\textasciitilde{}valtr/}{Research by P. Valtr was supported by the grant no. 18-19158S of the Czech Science Foundation (GA\v CR)}

\author{Uli Wagner}{IST Austria, Klosterneuburg, Austria}{uli@ist.ac.at}{https://orcid.org/0000-0002-3435-0100}{}

\authorrunning{R. Fulek and B. G{\"a}rtner and A. Kupavskii and P. Valtr and U. Wagner}%mandatory. First: Use abbreviated first/middle names. Second (only in severe cases): Use first author plus 'et al.'

\Copyright{Radoslav Fulek and Bernd G{\"a}rtner and Andrey Kupavskii and Pavel Valtr and Uli Wagner}%mandatory, please use full first names. LIPIcs license is "CC-BY";  http://creativecommons.org/li censes/by/3.0/

\ccsdesc[100]{Mathematics of computing~Combinatoric problems} \ccsdesc[100]{Theory of computation~Computational geometry}


% \ccsdesc[100]{Mathematics of computing~Combinatoric problems}, \ccsdesc[100]{Theory of computation~Computational geometry}
% mandatory: Please choose ACM 2012 classifications from https://www.acm.org/publications/class-2012 or https://dl.acm.org/ccs/ccs_flat.cfm . E.g., cite as "General and reference $\rightarrow$ General literature" or \ccsdesc[100]{General and reference~General literature}. 

\keywords{Discrete geometry, Tverberg theorem, Crossing Tverberg theorem}%mandatory

\category{}%optional, e.g. invited paper

\relatedversion{}%optional, e.g. full version hosted on arXiv, HAL, or other respository/website

\supplement{}%optional, e.g. related research data, source code, ... hosted on a repository like zenodo, figshare, GitHub, ...

\funding{}%optional, to capture a funding statement, which applies to all authors. Please enter author specific funding statements as fifth argument of the \author macro.

\acknowledgements{Part of the research leading to this paper was done during the 16th Gremo Workshop on Open Problems (GWOP), Waltensburg, Switzerland, June 12-16, 2018. We thank Patrick Schnider for suggesting the problem, and Stefan Felsner, Malte Milatz and Emo Welzl for fruitful discussions during the workshop. We also thank Stefan Felsner and Manfred Scheucher for finding and communicating the example from Section~\ref{sec:stronger}. We thank D{\"o}m{\"o}t{\"o}r P{\'a}lv{\"o}lgyi and the SoCG reviewers for various helpful comments.
}%optional
%Editor-only macros:: begin (do not touch as author)%%%%%%%%%%%%%%%%%%%%%%%%%%%%%%%%%%
\EventEditors{Gill Barequet and Yusu Wang}
\EventNoEds{2}
\EventLongTitle{35th International Symposium on Computational Geometry (SoCG 2019)}
\EventShortTitle{SoCG 2019}
\EventAcronym{SoCG}
\EventYear{2019}
\EventDate{June 18--21, 2019}
\EventLocation{Portland, United~States}
\EventLogo{socg-logo}
\SeriesVolume{129}
\ArticleNo{38} % <- Your article number goes here

\nolinenumbers %uncomment to disable line numbering
\hideLIPIcs  %uncomment to remove references to LIPIcs series (logo, DOI, ...), e.g. when preparing a pre-final version to be uploaded to arXiv or another public repository
%%%%%%%%%%%%%%%%%%%%%%%%%%%%%%%%%%%%%%%%%%%%%%%%%%%%%%
\begin{document}

\maketitle

\begin{abstract}
The Tverberg theorem is one of the cornerstones of discrete geometry. It states that, given a set $X$ of at least $(d+1)(r-1)+1$ points in $\mathbb R^d$, one can find a partition $X=X_1\cup \ldots \cup X_r$ of $X$, such that the convex hulls of the $X_i$, $i=1,\ldots,r$, all share a common point. In this paper, we prove a strengthening of this theorem that guarantees a partition which, in addition to the above, has the property that the boundaries of full-dimensional convex hulls have pairwise nonempty intersections. Possible generalizations and algorithmic aspects are also discussed. 

As a concrete application, we show that any $n$ points in the plane in general position span $\lfloor n/3\rfloor$ vertex-disjoint triangles that are pairwise crossing, meaning that their boundaries have pairwise nonempty intersections; this number is clearly best possible. A previous result of Alvarez-Rebollar et al.\ guarantees $\lfloor n/6\rfloor$ pairwise crossing triangles. Our result generalizes to a result about simplices in $\mathbb R^d,d\ge2$.
 \end{abstract}

\section{Introduction and Results}
The following theorem was published by Johann Radon in 1921 \cite{Radon1921}: any set of $d+2$ points in $\mathbb R^d$ can be partitioned into two (disjoint) subsets, whose convex hulls intersect. In 1966, Helge Tverberg \cite{T1966_tverberg} proved the following important generalization of Radon's result.
\begin{theorem}[Tverberg~\cite{T1966_tverberg}]
\label{thm:Tverberg}
Let $X$ be a set of at least $(d+1)(r-1)+1$ points in $d$-space. Then $X$ can be partitioned into $r$ sets whose convex hulls all have a point $o$ in common.
(In the literature, the point $o$ is referred to as a \emph{Tverberg point} and the partition as a \emph{Tverberg partition}.)
\end{theorem}
Radon's theorem covers the case $r=2$.
Figure~\ref{fig:tverberg_partition} illustrates the Tverberg theorem for $d=2$ and $r=3$.
\begin{figure}[htb]
\includegraphics[width=\textwidth]{tverberg_partition}
\caption{The Tverberg theorem for $d=2$ and $r=3$: any set of at least $7$ points can be partitioned into three sets whose convex hulls all have a point in common. Our example uses $9$ points and shows two Tverberg partitions as well as corresponding Tverberg points. Many other Tverberg partitions exist in this example.\label{fig:tverberg_partition}}
\end{figure}

This theorem largely influenced the course of discrete geometry and spurred a lot of research in the area. We do not go into more details in this paper and refer the reader to a recent survey by B{\'{a}}r{\'{a}}ny and Sober{\'{o}}n~\cite{Barany2018}.

Throughout this paper, a set of points in $\R^d$ is said to be in \emph{general position} if no $d+1$ points are on a comon hyperplane (for example, no three points on a line in the plane).  

There is a major open question from the field of geometric graphs, motivating our work. In \cite{Aronov1994}, Aronov et al.\  conjectured that there exists an absolute constant $c>0$, such that, given any set of $n$ points in general position in the plane, one can find at least $cn$ disjoint pairs among them such that their connecting segments cross pairwise. Such a collection of segments is called a \emph{crossing family}. Despite considerable interest in this problem, the best published bound still comes from the original paper \cite{Aronov1994}, stating that one can always find a crossing family of size at least $c\sqrt n$ for some absolute $c>0$.
Only recently, after about 25 years, J.~Pach, N.~Rubin, and G.~Tardos improved this lower bound to a nearly linear one~\cite{PRT19_crossing}.

In another attempt to approach this problem, Alvarez-Rebollar et al.\ asked whether one can find at least $cn$ disjoint \emph{triples} whose connecting \emph{triangles} cross pairwise~\cite{RLU2016_crossingFam}. Throughout this paper, we say that two convex polytopes in $\R^d$ \emph{cross} if their boundaries have a non-empty intersection. We remark that if a convex polytope in $\R^d$ is not full-dimensional then it has no interior and therefore coincides with its boundary. Our main results below would be false, in the absence of general position, if relative boundaries were considered in the above definition of crossing (an easy counterexample in this situation is a set of points lying on a line in $\R^2$). 

We also point out that our definition of ``crossing'' is different from a classical one due to Fejes-T\'oth for general convex sets~\cite{FT67}. According to the latter, two convex sets $C,C'$ cross if neither $C\setminus C'$ nor $C'\setminus C$ are connected. Under general position, our notion of crossing convex polytopes generalizes the notion of a crossing of two vertex-disjoint segments in a crossing family (the two segments intersecting in their relative interiors)~\cite{Aronov1994}. It also generalizes the notion of Alvarez-Rebollar et al.\ for vertex-disjoint triangles (an edge of one triangle crossing an edge of the other)~\cite{RLU2016_crossingFam}.

Alvarez-Rebollar et al.\ showed the following: For every finite point set of size $n$ in the plane in general position, there exist $\left\lfloor\frac{n}{6}\right\rfloor$ vertex-disjoint and pairwise crossing triangles with vertices in $P$. As at most $\left\lfloor\frac{n}{3}\right\rfloor$ disjoint triangles can be found, this leaves a factor-2 gap.

If we only want triangles that have a common point, this gap can be closed using a simple strengthening of the Tverberg theorem for point sets of size at most $(d+1)r$ that we present next. However, these triangles might not be pairwise crossing, since triangles can be nested; see Figure~\ref{fig:tverberg_partition} (right).

\begin{theorem}\label{thm:Tverberg2}
Let $X$ be a set of at least $(d+1)(r-1)+1$ and at most $(d+1)r$ points in $d$-space. Then $X$ can be partitioned into $r$ disjoint sets $X_1,\ldots, X_r$ of size at most $d+1$, whose convex hulls all have a point in common. 
\end{theorem}

To prove this, we apply Theorem~\ref{thm:Tverberg} to $X$. We get sets $X'_1,\ldots, X'_r$ whose convex hulls contain a common point, say, the origin. Using Carath\' eodory's theorem, from every $X'_i$ of size larger than $d+1$ we can select $d+1$ points $X_i\subseteq X'_i$, whose convex hull still contains the origin. Finally, some of the sets $X'_i$ of size smaller than $d+1$ are filled up to size $d+1$ with the points removed from other $X'_i$'s.  The origin is still a common point for all $\conv(X_i)$'s.

The main result of this paper provides a crossing version of the Tverberg Theorem~\ref{thm:Tverberg2}.
%, also closing the gap in the result of Rebollar et al.~\cite{RLU2016_crossingFam}, and generalizing it to higher dimensions.

\begin{theorem} \label{thm:main} Let $X$ be a set of at least $(d+1)(r-1)+1$ and at most $(d+1)r$ points in $d$-space. Then $X$ can be partitioned into $r$ disjoint sets $X_1,\ldots, X_r$ of size at most $d+1$, whose convex hulls all have a point in common. Moreover, for any $X_i,X_j$ of size $d+1$, $\conv (X_i)$ and $\conv (X_j)$ cross, meaning that their boundaries have a non-empty intersection.
\end{theorem} 

We call such a partition a \emph{crossing Tverberg partition}. An easy calculation shows that the number of sets with exactly $d+1$ elements is at least $|X|-dr\in\{r-d,r-d+1,\ldots,r\}$. In particular, for sets $X$ of size exactly $(d+1)r$, we immediately deduce the following simpler-looking corollary.

\begin{corollary}\label{cor1} Let $X$ be a set of $(d+1)r$ points in $d$-space. Then $X$ can be partitioned into $r$ sets $X_1,\ldots, X_r$ of size $d+1$, whose convex hulls all have a point in common and such that for any $i,j\in [n]$, $\conv (X_i)$ and $\conv (X_j)$ cross, meaning that their boundaries have a non-empty intersection.  
\end{corollary}

We also obtain an optimal strenghtening of the result by Alvarez-Rebollar et al.~\cite{RLU2016_crossingFam} that moreover generalizes to all dimensions $d\geq 2$.

\begin{corollary}
For every finite point set $X$ of size $n$ in the plane in general position, there exist $\lfloor n/3\rfloor$ vertex-disjoint and pairwise crossing triangles with vertices in $X$.

More generally, for every finite point set $X$ of size $n$ in $\mathbb R^d$ in general position, there exist $\lfloor n/(d+1)\rfloor$ vertex-disjoint and pairwise crossing simplices with vertices in $X$.
\end{corollary}
To derive this from Theorem~\ref{thm:main}, we remove $n\mod (d+1)$ points from $X$ and then apply Theorem~\ref{thm:main} on the remaining set of $(d+1)\left\lfloor\frac{n}{d+1}\right\rfloor$ points.

Finally, we get a crossing version of the actual Tverberg Theorem~\ref{thm:Tverberg}, for point sets of arbitrarily large size.

\begin{theorem}[Crossing Tverberg theorem]\label{thm:crossing_tverberg}
Let $X$ be a set of at least $(d+1)(r-1)+1$ points in $d$-space. Then $X$ can be partitioned into $r$ sets whose convex hulls all have a point in common. Moreover, for any $X_i,X_j$ of size at least $d+1$, $\conv (X_i)$ and $\conv (X_j)$ cross, meaning that their boundaries have a non-empty intersection.  
\end{theorem} 

This is also easy to prove, using Theorem~\ref{thm:main} and Corollary~\ref{cor1}. If $n:=|X|\leq (d+1)r$, we apply Theorem~\ref{thm:main}. Otherwise, we apply Corollary~\ref{cor1} to an arbitrary subset $Y\subset X$ of size $(d+1)r$, resulting in a crossing Tverberg partition into $r$ sets $Y_1,\ldots,Y_r$ of size $d+1$ each. Now we consecutively add the remaining points to suitable sets in such a way that the crossings between the convex hulls are maintained (the Tverberg point automatically remains valid). Suppose that some points have already been added, resulting in sets $X'_1,\ldots,X'_r$ whose convex hulls still cross pairwise. If the next point $p$ is contained in one of these convex hulls, we simply add it to the corresponding set. Otherwise, we select an inclusion-minimal set $\conv(X'_k\cup\{p\})$ among the (different) sets $\conv(X'_i\cup\{p\}), i=1,\ldots,r$, and add $p$ to $X'_k$.  We claim that this cannot result in nested convex hulls. For this, we only have to rule out that $\conv(X'_k\cup\{p\})$ ``swallows'' some other $\conv(X'_i)$. But this cannot happen, as otherwise $\conv(X'_i\cup\{p\})\subset \conv(X'_k\cup\{p\})$, contradicting the choice of $k$.

It remains to prove Theorem~\ref{thm:main}. 
We start with a Tverberg partition according to Theorem~\ref{thm:Tverberg2}, i.e.\ a partition into sets of size at most $d+1$ such that their convex hulls have a point in common. Such a partition might look like in Figure~\ref{fig:tverberg_partition2} (left). If the full-dimensional convex hulls (which are simplices) cross pairwise, we are done. In general, however, we still have pairs of nested simplices. As long as this is the case, we \emph{fix} one pair of nested simplices at a time until no pair of nested simplices exists anymore, and our desired crossing Tverberg partition is obtained.

By fixing, we mean that we repartition the $2(d+1)$ points involved in the two simplices in such a way that the resulting simplices are not nested anymore but still contain the common point. In the example of Figure~\ref{fig:tverberg_partition2} (left), there is one pair of nested simplices (red edges), and after fixing it (blue edges), we are actually done in this case; see Figure~\ref{fig:tverberg_partition2} (right).

\begin{figure}[htb]
\includegraphics[width=\textwidth]{tverberg_partition2}
\caption{Fixing a pair of nested simplices spanned by the 6 black points: two nested simplices (red) are transformed into two simplices that cross (blue).\label{fig:tverberg_partition2}}
\end{figure}

% There are of course some things to prove here. First of all, we need to show that we can actually perform the desired repartitioning in all cases. Then, by fixing one pair, new pairs might arise that need fixing, so it has to be proved that fixing terminates.

Repartitioning the points is always possible due to the following lemma, which may be of independent interest. This lemma is our main technical contribution, and we prove it in Section~\ref{sec:parity}.

\begin{lemma}\label{lem:unnesting}
Let $T,T'$ be two disjoint $(d+1)$-element sets in $\mathbb R^d$ such that $\nula \in \conv(T)\cap\conv(T')$. Then there exist two disjoint $(d+1)$-element sets $S,S'$ such that $S\cup S'=T\cup T'$, $\nula \in \conv(S)\cap\conv(S')$, and moreover, $\conv(S)$ and $\conv(S')$ cross.
\end{lemma}

To conclude the proof of Theorem~\ref{thm:main}, we still need to show that after finitely many fixing steps, there are no nested pairs anymore, in which case we have a crossing Tverberg partition.

To see this, we observe that in any fixing operation that replaces $X_i$ and $X_j$ such that $\conv(X_i)\subset \conv(X_j)$, the simplex $\conv(X_j)$ is volume-wise the unique largest $d$-dimensional simplex that can be formed from the $2(d+1)$ points $X_i\cup X_j$ involved in the operation. Hence, the two simplices $\conv(S)$ and $\conv(S')$ replacing $\conv(X_i)$ and $\conv(X_j)$ are volume-wise both strictly smaller than $\conv(X_j)$. Therefore, if we order all full-dimensional simplices by decreasing volume, the sequence of these volumes goes down lexicographically in every fixing operation.

Formally, let ${\cal V}=(V_1,V_2\ldots V_s), s\leq r$ be the sequence of volumes in decreasing order before the fix, and ${\cal V}'=(V'_1,V'_2\ldots V'_s)$ the decreasing order after the fix. Moreover, suppose that $k$ is the largest index at which the volume of $\conv(X_j)$ appears in ${\cal V}$. 
% (by even more general position, we could assume that there is a unique such index, but this does not really help here). 
The volume of $\conv(X_i)$ appears at a position $\ell>k$. As the fixing operation removes a volume equal to $V_k$ and inserts two volumes smaller than $V_k$, we have that ${\cal V}$ and ${\cal V}'$ agree in the first $k-1$ positions, but $V'_{k}<V_k$. This exactly defines the relation ``${\cal V'}<{\cal V}$'' in decreasing lexicographical order, and as this is a total order, fixing must eventually terminate.

\vskip+0.1cm
\textbf{Remark. } Instead of volume, we may use the number of points from $X$ inside $\conv(X_j)$ as a measure. 

Since applications of Lemma~\ref{lem:unnesting} allow to keep the Tverberg point and the sizes of the parts in the (Tverberg) partition, we actually have the following strengthening of Theorem~\ref{thm:main}.

\begin{theorem}\label{thm:main-stronger} 
Let $X$ be a set of points in $d$-space.  Suppose that there is a Tverberg partition of $X$ into $r$ parts of sizes $s_1\dots,s_r$, for which $o$ is a Tverberg point, and $s_i\leq d+1, i=1,\ldots r$. Then there is also a \emph{crossing} Tverberg partition of $X$ into $r$ parts of sizes $s_1\dots,s_r$, for which a Tverberg point is $o$.  
\end{theorem}

In the next section, we prove Lemma~\ref{lem:unnesting}. We remark that for each of the other theorems and corollaries in this section, we have already explained how to derive them from the Tverberg theorem, or from our main Theorem~\ref{thm:main}. In Section~\ref{sec3}, we discuss possible generalizations as well as some algorithmic aspects of the problem.
%In fact, we can prove something stronger: not only the aforementioned $d$-simplices mutually intersect, but they actually all share a  common point.

%In the following, we slightly abuse notation and identify a (d+1)-element subset $T\subseteq P$ with the simplex it defines (the convex hull of $T$).

% \section{Proof of Theorem~\ref{thm:main}}

% % First we show that we can make the following assumptions: (A1) $X$ is in general position, meaning that no $d+1$ points of $X$ lie in a common hyperplane; equivalently, every subset $S\subseteq X$ of size $d+1$ spans a full-dimensional simplex $\conv(S)$. (A2) For any collection ${\cal S}\subseteq\binom{X}{d+1}$ of disjoint $d+1$-element subsets of $X$, the intersection $\cap_{S\in{\cal S}} \conv(S)$ of the spanned simplices is either empty, or full-dimensional.

% % To justify these assumptions, we make two observations. The first one is that (A1) and (A2) hold for \emph{some} arbitrarily small (for example, random) perturbation $X'$ of $X$. The second is that the property of not having a crossing Tverberg partition is maintained under \emph{any} sufficiently small perturbation. Putting this together, we see that if every point set $X'$ satisfying (A1) and (A2) has a crossing Tverberg partition, then every point set $X$ has a crossing Tverberg partition. 

% %Henceforth, we assume (A1) and (A2) without explicitly mentioning it.

% % union of the $X_i$'s of size $d+1$ and the common point of all convex hulls are in general position. To justify assumption (A1), we observe that if $X'$ is a sufficiently small perturbation of $X$ then every crossing Tverberg partition of $X'$ corresponds to a crossing Tverberg partition of $X$.
% % Assumption (A2) is then satisfied for any (Tverberg) point which lies in the relative interior of the set of all Tverberg points for all the sets $X_i$, as otherwise we would get a contradiction with assumption (A1).


% %First of all, using a standard topological argument, we may assume that the points of $X$ are in general position, meaning that no $d+1$ points lie in a common hyperplane. To justify this assumption, we observe that the set of point sets of some fixed size $n$ that allow a crossing Tverberg partition is closed, since its complement is open. Indeed, if we take a point set that does not allow a crossing Tverberg partition, then all (finitely many) partitions have the property that there is no common point in all the convex hulls, or that some convex hulls are (properly) nested. This property is maintained under any sufficiently small perturbation. Since every point set is the limit of a sequence of point sets in general position, it follows from closedness that if all points sets in general position allow a crossing Tverberg partition, then all point sets do. 

% %Under this general position assumption, we may also assume that the union of the $X_i$'s of size $d+1$ and the common point of all convex hulls are in general position. Henceforth, we assume general position without explicitly mentioning it.

% We start with a Tverberg partition of $X$ into sets $X_1,\ldots,X_r$ of size at most $d+1$; such a partition is guaranteed by the Tverberg Theorem~\ref{thm:Tverberg2}. 

% % Under (A2), we can slightly perturb the Tverberg point $o$ such that---under (A1)---not only $X$, but $X\cup\{o\}$ is in general position. Finally, we assume w.l.o.g.\ that the Tverberg point is the origin, $o=\nula$.

% \subsection{\emph{Fixing Pairs}}
% The sets $X_1,\ldots,X_{r}$ may not yet satisfy the requirements of Theorem~\ref{thm:main}, since there may be pairs $\{X_i, X_j\}$, $X_i\not= X_j$ such that $\conv(X_i)$ and $\conv(X_j)$ do not cross. As the simplices $\conv(X_i)$ and $\conv(X_j)$ themselves intersect in the origin, the only remaining possibility is that one of them properly contains the other one, i.e., they are nested; see Figure~\ref{fig:tverberg_partition} (right). As long as there exists a pair of nested simplices $\conv(X_i)\subset\conv(X_j)$, we apply the \emph{fixing pairs} operation, i.e.\ we use Lemma~\ref{lem:unnesting} to replace $T=X_i$ and $T'=X_j$ with $S,S'$ such that $\conv(S)$ and $\conv(S')$ cross. 

% \subsection{\emph{Fixing Pairs} terminates}
% We need to show that after finitely many steps, there are no nested pairs anymore, in which case we have a crossing Tverberg partition.

% To see termination, we observe that in any \emph{fixing pairs} operation that replaces $X_i$ and $X_j$ such that $\conv(X_i)\subset \conv(X_j)$, the simplex $\conv(X_j)$ is volume-wise the unique largest $d$-dimensional simplex that can be formed from the $2(d+1)$ points $X_i\cup X_j$ involved in the operation. Hence, the two simplices $\conv(S)$ and $\conv(S')$ replacing $\conv(X_i)$ and $\conv(X_j)$ are volume-wise both strictly smaller than $\conv(X_j)$. Therefore, if we order all full-dimensional simplices by decreasing volume, the sequence of these volumes goes down lexicographically in every fixing operation.

% Formally, let ${\cal V}=(V_1,V_2\ldots V_s), s\leq r$ be the sequence of volumes in decreasing order before the fix, and ${\cal V}'=(V'_1,V'_2\ldots V'_s)$ the decreasing order after the fix. Moreover, suppose that $k$ is the largest index at which the volume of $\conv(X_j)$ appears in ${\cal V}$. 
% % (by even more general position, we could assume that there is a unique such index, but this does not really help here). 
% The volume of $\conv(X_i)$ appears at a position $\ell>k$. As the fixing operation removes a volume equal to $V_k$ and inserts two volumes smaller than $V_k$, we have that ${\cal V}$ and ${\cal V}'$ agree in the first $k-1$ positions, but $V'_{k}<V_k$. This exactly defines the relation ``${\cal V'}<{\cal V}$'' in decreasing lexicographical order, and as this is a total order, \emph{fixing pairs} must eventually terminate.



\section{Proof of Lemma~\ref{lem:unnesting}}\label{sec:parity}
We employ the following parity lemma that we prove in the next two subsections (in a geometric way for $d=2$, and in a combinatorial way for all dimensions).

\begin{lemma}[Parity lemma]\label{lemparity} Let $V\subset \mathbb R^d$, $|V|=2(d+1)$, such that $V\cup\{\nula\}$ is in general position. Then
$$\Big|\Big\{\{F,G\}:F,G\in {V\choose d+1}, F\cap G=\emptyset, \nula\in \conv F\cap \conv G\Big\}\Big| \ \ \ \text{is even.}$$
\end{lemma}
In particular, if there is one such pair $\{F,G\}$, then there is another.

Before we move on to the proof, let us derive Lemma~\ref{lem:unnesting} from this. It suffices to handle the case where the set $T\cup T'\cup\{\nula\}$ is in general position. To justify this, we make two observations. The first is that general position is achieved by \emph{some} arbitrarily small (for example, random) perturbation of $T\cup T'$. The second is that the property of \emph{not} having sets $S,S'$ as required is maintained under \emph{any} sufficiently small perturbation. Putting the two observations together, we see that if Lemma~\ref{lem:unnesting} holds for $T\cup T'\cup\{\nula\}$ in general position, then it holds for all $T,T'$.

Then, by Lemma~\ref{lemparity} with $V=T\cup T'$, there is another pair $\{S,S'\}$ such that $S\cup S'=T\cup T'$, $\nula \in \conv(S)\cap\conv(S')$. At most one of these pairs can yield nested simplices: the outer simplex of a nested pair coincides with $\conv(T\cup T')$ and is therefore uniquely determined. Hence, one of the two pairs yields crossing simplices. 

\vskip+0.1cm
\textbf{Remark. } We note that this statement and its application is quite unusual. Typically, one proves that the number of objects of a certain type is always odd and thus at least one object of the type exists.

\subsection{Geometric proof of the parity Lemma~\ref{lemparity} for $\mathbf{d=2}$}
The purely combinatorial proof that we give below for all $d$ is not difficult, but does not provide any geometric intuition. We therefore start with a simple proof in the plane.

Consider a set $V$ of $2(d+1)=6$ points in the plane. We remark that the statement is invariant under scaling points, and thus we may assume that all points lie on a circle with the center in the origin. Due to general position, no two points from $V$ lie on a line passing through the origin; see Figure~\ref{fig:geomproof} (left).

\begin{figure}[htb]
\begin{center}
\includegraphics[width=0.6\textwidth]{geomproof}
\caption{6 points in the plane (black) and their mirror images (white); there must be two consecutive black points $p,p'$.\label{fig:geomproof}}
\end{center}
\end{figure}

Now we reflect each point at the origin and obtain another $6$ points, drawn in white in  Figure~\ref{fig:geomproof} (right). We observe that in the circular order of points, there must be two consecutive ones of the same color. Indeed, an alternating pattern (black, white, black, white,\ldots) would lead to pairs of antipodal points having the \emph{same} color. Let $p,p'$ be two consecutive points of the same color; by going to the antipodal points if necessary, we may assume that they are black and hence belong to $V$. 

We make two observations (actually, just one). (i) $p$ and $p'$ cannot belong to a triple $F\subset V$ such that $\nula\in\conv(F)$, as otherwise, the third point would get reflected to a white point between $p$ and $p'$; see Figure~\ref{fig:geomproof2} (left). (ii) For any $q\in V\setminus\{p,p'\}$, the two segments $\conv(\{p,q\})$ and $\conv(\{p',q\})$ have the origin on the same side; see Figure~\ref{fig:geomproof2} (right). 

\begin{figure}[htb]
\begin{center}
\includegraphics[width=0.6\textwidth]{geomproof2}
\caption{Consecutive black points are combinatorially indistinguishable.\label{fig:geomproof2}}
\end{center}
\end{figure}

This implies the following: if $\{F,G\}$ is a partition of $V$ that we count in Lemma~\ref{lemparity}, then (i) $p$ and $p'$ are in different parts, and (ii) swapping $p$ and $p'$ between the parts leads to a different partition $\{F',G'\}$ that we also count. In other words, $p$ and $p'$ are ``combinatorially indistinguishable'' with respect to the relevant properties, and the operation of swapping them between parts establishes a matching between the partitions that we want to count. Hence, their number is even.

% In this section, we sketch a geometric proof for Lemma~\ref{lemparity} for $d=2$. It easily extends to higher dimensions. However, it relies to some extent on geometric intuition, and thus in the next section we give a more formal combinatorial proof. One may actually see that the steps of this proof have their combinatorial counterparts in the proof from the next section.

% Consider a set $V$ of $6$ points on the plane and triangles $T,T'$ as in Lemma~\ref{lem:unnesting}. We remark that the problem is actually invariant under multiplying points by scalars, and thus we may assume that all points lie on a unit circle with the center in the origin. Due to our general position assumption, no two points from $V$ lie on a line passing through the origin.

% We call any partition of $V$ into two vertex-disjoint triangles a {\it split}. If additionally both triangles contain the origin, then we call it a {\it good split}. We shall prove that, for any set $V$ of $6$ points on the circle in general position, there is an even number of good splits of $V$. We start with a configuration in which all points, say, have positive $x$-coordinate and thus there are no good splits. Next, we move points one by one to their position in $V$.

% The number of good splits changes only when the point $p$ being moved passes through a degeneracy. We call a configuration {\it degenerate} if both $p$ and $-p$ are in the set. Moreover, for a split to become/stop being good after $p$ passes through a degeneracy, the vertices of the triangles in the split must include both $p$ and $-p$. For simplicity, we assume in what follows that $p=(0,1)$ and that $p$ moves in the direction of increasing the $x$-coordinate.

% Look at the degenerate configuration and the set $V':=V\setminus \{p,-p\}$. $V'$ consists of $4$ points, and thus may be seen as a projection of some tetrahedron. The line in $\mathbb R^3$ that projects into the origin intersects either $2$ or $0$ faces of this tetrahedron. Thus, there are either $2$ or $0$ triangles with vertices in $T$ that contain the origin. If there are no such triangles, then there are no good splits, in which one triangle has vertices $p$ and $-p$, neither before nor after the degeneracy. 

% Assume next that there are two triangles $T_1,T_2\subset V'$ that contain the origin. They share a common edge $E$. Let $p_i = T_i\setminus E$, $i=1,2$. Due to the discussion above, any split that becomes/stops being good after passing through the degeneracy must have the form $\{\{p_i\}\cup E, \{p_{i'},p,-p\}\}$, where $\{i,i'\}=\{1,2\}$. Given the form of $p$, we arrive to the following conclusion:
% \begin{itemize}
%   \item If both $p_1,p_2$ have (strictly) negative $x$-coordinate then both splits above were bad and became good after passing through the degeneracy. Thus, the number of good splits increases by $2$.
%   \item If  one of $p_1,p_2$ has (strictly) negative $x$-coordinate and the other one has (strictly) positive coordinate, then the split $\{\{p_1\}\cup E, \{p_{2},p,-p\}\}$ was good and it stops being good after passing through degeneracy. The situation with the other split is opposite, and thus the number of good splits stays the same.
%   \item If both $p_1,p_2$ have (strictly) positive $x$-coordinate then both splits above were good and became bad. Thus, the number of good splits decreases by $2$.
% \end{itemize}
% In any case, the number of good splits stays even. This concludes the proof of the lemma.

%%% FIX STARTS HERE %%% 
%%% --------------  %%%
\subsection{Combinatorial proof of the parity Lemma~\ref{lemparity}}
Recall that for $|V|=2(d+1)$, we need to show that there is an even number of partitions $\{F,G\}$ of $V$ into parts of equal size $d+1$ such that $\nula\in\conv(F)\cap\conv(G)$.
We show that this follows from the fact that the $(d+1)$-element subsets $F$ with $\nula\in\conv(F)$ form a \emph{cocycle}, a concept borrowed from topology. %The proof itself does not use any topology, though. 

\begin{definition} Let $n\ge k\ge 1$ be integers, and let $V$ be a set with $n$ elements. A family $\mathcal C\subset {V\choose k}$ of $k$-element subsets of $V$ is a {\rm cocycle} if $$|\{F\in \mathcal C: F\subset M\}| \ \ \text{is even for every } M\subset V\ \ \text{with } \ |M|=k+1. $$\end{definition}

\textbf{Example. } Fix a set $D\in {V\choose k-1}$. Then $\delta D:=\{F\in {V\choose k}: D\subset F\}$ is a cocycle. Indeed, a $(k+1)$-element subset $M\subset V$ either contains no sets in $\delta D$ (if $D\not \subset M$) or exactly two sets in $\delta D$ (If $D=M\setminus\{p,q\}$).

\begin{lemma}\label{lem1}
  Let $V\subset \mathbb R^d$ be such that $V\cup\{\nula\}$ is in general position. Then
  $$\mathcal C(V):=\big\{F\in {V\choose d+1}: \nula \in \conv F\big\}$$
  is a cocycle.
\end{lemma}
\begin{proof}
Let $M:=\{v_1,\ldots,v_{d+2}\}\subset V$. Lift the points to dimension $d+1$ such that the convex hull of the lifted point set $\hat{M}$ is a full-dimensional simplex $\Delta$; see Figure~\ref{fig:lifting}.

\begin{figure}[htb]
\begin{center}
\includegraphics[width=0.6\textwidth]{lifting}
\caption{Illustration of the geometric proof of Lemma~\ref{lem1} when $d=2$\label{fig:lifting}}
\end{center}
\end{figure}

Any set $\hat{F}\subset \hat{M}$ of size $d+1$ spans a facet of $\Delta$, and we have $\nula\in\conv(F)$ if and only if the vertical line through $\nula$ intersects that facet. As $V\cup\{\nula\}$ is in general position, this line intersects the convex set $\Delta$ in a line segment, if it intersects it at all; hence, the line intersects the boundary of $\Delta$ in zero or two points, each lying in the interior of some facet. Thus, $|\{F\in \mathcal C(V): F\subset M\}|\in\{0,2\}$.

We remark that there is also an elementary linear algebra version of this ``proof by picture'' (omitted in this extended abstract). 
\end{proof}

% Consider the $(d+1)\times (d+2)$-matrix 
% \[
% A = \left(
% \begin{array}{ccc}
% v_1 & \cdots & v_{d+2} \\
% 1 & \cdots & 1 
% \end{array}
% \right).
% \]
% By general position of $V$, every $(d+1)\times (d+1)$ submatrix of $A$ is invertible, and in particular, $A$ itself has full rank $d+1$. Indeed, a linear dependence between the rows of a $(d+1)\times (d+1)$ submatrix would induce a hyperplane containing $d+1$ of the points from $V$. It follows that for $F=M\setminus\{v_k\}$, the system of equations
% \[
% \sum_{i\neq k}\lambda_i v_i = \nula, \quad \sum_{i\neq k}\lambda_i=1 
% \]
% along with $\lambda_k=0$ has a unique solution $\boldsymbol{\lambda}^F\in\R^{d+2}$. Moreover, $\boldsymbol{\lambda}^F_i\neq 0$ for all $i\neq k$, as $\nula$ cannot be an affine combination of less than $d+1$ points, by general position of $V\cup\{\nula\}$. We have 
% $\boldsymbol{\lambda}^F\geq 0$ if and only if $F\in\mathcal C$, hence we need to show that there is an even number of such $\boldsymbol{\lambda}^F$. 
% If $\nula\notin\conv(M)$, then there is none. Otherwise, 
% Carath\' eodory's Theorem guarantees the existence of some $\boldsymbol{\lambda}^G\geq 0$, and we need to show that there is exactly one more $\boldsymbol{\lambda}^H\geq 0$. To see this, we observe that all $\boldsymbol{\lambda}^F$ are on the line $\{\boldsymbol{\lambda}^G+t\boldsymbol{\lambda}: t\in\R\}$, where $\boldsymbol{\lambda}$ is some nonzero solution of the system
% \[
% \sum_{i=1}^{d+2}\lambda_i v_i = \nula, \quad \sum_{i=1}^{d+2}\lambda_i=0. 
% \]
% Since matrix $A$ has rank $d+1$, the solution space of this system is one-dimensional. We also have that $\boldsymbol{\lambda}_i\neq 0$ for all $i$, as no $d+1$ columns of $A$ are linearly dependent. Moreover, because of the equation $\sum_{i=1}^{d+2}\lambda_i=0$, there are positive as well as negative entries in $\boldsymbol{\lambda}$. Suppose that $G=M\setminus\{v_{\ell}\}$ and assume w.l.o.g.\ that $\boldsymbol{\lambda}_{\ell}>0$. Then, any $\boldsymbol{\lambda}^F\geq 0$ is of the form $\boldsymbol{\lambda}^G+t\boldsymbol{\lambda}$ for $t\geq 0$. For $t=0$, we get $\boldsymbol{\lambda}^G$. We let $t$ increase until some entry of $\boldsymbol{\lambda}^G+t\boldsymbol{\lambda}$ becomes zero. As $\boldsymbol{\lambda}^G_i>0$ for $i\neq\ell$ and $\boldsymbol{\lambda}$ has a negative entry, this happens at some positive value $t^*$, and $\boldsymbol{\lambda}^H:=\boldsymbol{\lambda}^G+t^*\boldsymbol{\lambda}\geq 0$ is our second nonnegative $\boldsymbol{\lambda}^F$. For $t>t^*$, we have at least one negative component, so there are no further nonnegative $\boldsymbol{\lambda}^F$.


% Consider the standard $(d+1)$-dimensional simplex $\Delta\subset \R^{d+1}$ with vertices $e_1,\ldots, e_{d+2}$, and consider the affine map $A: \R^{d+1}\to \R^d$ given by $e_i\mapsto v_i,$ $i=1,\ldots, d+2$.

%   By the general position assumption, the preimage $A^{-1}(\nula)$ is a line in general position with respect to the simplex $\Delta$. Thus, the line $A^{-1}(\nula)$ intersects either zero or two facets of $\Delta$.


By Lemma~\ref{lem1}, Lemma~\ref{lemparity} is now simply a special case of the following main result of this section.

\begin{theorem}\label{prop2} Let $k\ge 1$, $|V|=2k$ and let $\mathcal C\subset {V\choose k}$ be a cocycle. Set    $$\mathcal P_{\mathcal C}:=\big\{\{F,G\}: F,G\in \mathcal C, F\cap G=\emptyset \big\}.$$
Then $|\mathcal P_{\mathcal C}|$ is even.
\end{theorem}


For the proof, we need one more concept. For a family $\mathcal D\subset {V\choose k-1}$, we define the \emph{coboundary} 
\[\delta\mathcal D=\big\{F\in {V\choose k}: F\ \text{contains an odd number of sets }D\in \mathcal D\big\}.\] This naturally extends the definition $\delta D=\{F\in {V\choose k}: D\subset F\}$ for $D\in {V\choose k-1}$ from our previous example. We call $\delta D$ the \emph{elementary cocycle} associated with $D$. 

To prove Theorem~\ref{prop2}, we use the following basic fact.

\begin{lemma}\label{lem2}
For every cocycle $\mathcal C\subset {V\choose k}$ there exists a family $\mathcal D\subset {V\choose k-1}$ such that
$\mathcal C=\delta\mathcal D$.
Equivalently,
$$\mathcal C=\delta D_1\oplus\delta D_2\oplus \ldots \oplus \delta D_m,$$
where $\mathcal D=\{D_1,\ldots, D_m\}$ and $\oplus$ denotes symmetric difference.  
\end{lemma}

\begin{proof}[Proof of Theorem~\ref{prop2}]
The statement is trivially true when $\mathcal C = \emptyset$ is the empty cocycle.

Thus, by Lemma~\ref{lem2}, it suffices to show that the parity of $|\mathcal P_{\mathcal C}|$ does not change when we take the symmetric difference with an elementary cocycle, i.e., 
\begin{equation}\label{eqparity}|\mathcal P_{\mathcal C}|\equiv |\mathcal P_{\mathcal C\oplus\delta D}|\mod 2\end{equation}
for any $D\in {V\choose k-1}$.
\begin{claim} $\mathcal P_{\mathcal C\oplus\delta D}=\mathcal P_{\mathcal C}\oplus \mathcal R,$ where
$$\mathcal R:=\big\{\{F,G\}: F\in \delta D, G\in \mathcal C, F\cap G=\emptyset \big\}.$$
\end{claim}
\begin{proof}
To see this, note that for pairs $\{F,G\}$ with $D\not \subset F$ and $D\not\subset G$ nothing changes, i.e., such a pair is either in both $\mathcal P_{\mathcal C}$ and $\mathcal P_{\mathcal C\oplus\delta D}$ or in neither. Thus, it suffices to consider pairs $\{F,G\}$, $F\cap G=\emptyset$, such that $D$ is contained in one of the two sets. Without loss of generality, $D\subset F$ and $G\subset V\setminus D$.

Since $D\not\subset G$, we have $G\in \mathcal C$ if and only if $G\in \mathcal C\oplus\delta D$.  Since $D\subset F$, we have $F\in \delta D$, and hence, $F\in \mathcal C$ if and only if $F\not\in \mathcal C\oplus \delta D$. Thus, $\{F,G\} \in \mathcal P_{\mathcal C\oplus\delta D}$ if and only if  $\{F,G\}\in \mathcal{R}$ and  $\{F,G\}\not\in \mathcal{P}_{\mathcal{C}}$. This proves the claim.
%\vskip+0.1cm
\end{proof}

Since $\mathcal C$ is a cocycle and $|V\setminus D|=k+1$, there is an even number of $G\in \mathcal C$ with $G\subset V\setminus D$. Since $|V|=2k$, each such $G$ determines a unique $F=V\setminus G$, $F\in \delta D$, with $\{F,G\}\in \mathcal R$. Thus,
$$|\mathcal R|=\big|\{G\in \mathcal C: G\subset V\setminus D\}\big|\equiv 0\mod 2.$$
Combined with the claim, we get that \eqref{eqparity} holds, which concludes the proof of the theorem.
\end{proof}

For the sake completeness, we also prove Lemma~\ref{lem2}.

\begin{proof}[Proof of Lemma~\ref{lem2}] Let $\mathcal C\subset {V\choose k}$ be a cocycle. Choose an arbitrary element $v\in V$. Define 
$$\mathcal D:=\mathcal D_v:=\big\{D\in{V\setminus v\choose k-1}: D\cup \{v\}\in \mathcal C\big\}.$$
 
We claim that $$\mathcal C=\delta \mathcal D.$$%:=\big\{F\in {V\choose k-1}: F\text{ contains an odd number of sets }D\in\mathcal D\big\}.$$
Let $F\in {V\choose k}$. We distinguish two cases:
\begin{itemize}
  \item[1) ] If $v\in F$ then, by the definition of $\mathcal D\subset {V\setminus \{v\}\choose k-1}$, we have $F\setminus \{v\}\in \mathcal D$ if and only if $F\in \mathcal C$. Moreover, $F\setminus \{v\}$ is the only set in ${V\setminus \{v\}\choose k-1}$ that $F$ contains. Thus,
$$F\in \mathcal C\ \Leftrightarrow\ F\in \delta \mathcal D\ \ \ \text{in this case.}$$
  \item[2) ] Assume $v\notin F$. Then $M:=F\cup \{v\}$ has $k+1$ elements. Since $\mathcal C$ is a cocycle, $M$ contains an even number of sets from $\mathcal C$. Thus,
\begin{align*}
  F\in \mathcal C & \Leftrightarrow M \ \text{ contains an odd number of }G\in \mathcal C \ \text{with } v\in G \\
   & \Leftrightarrow F\ \text{ contains an odd number of sets }D=G\setminus \{v\}\in \mathcal D\\
   & \Leftrightarrow F\in \delta \mathcal D.
\end{align*}
\end{itemize}
\end{proof}
%%% FIX ENDS HERE %%% 
%%% ------------- %%%

\section{Discussion}\label{sec3}
\subsection{A Topological version of Theorem~\ref{thm:main}}

%For appropriately chosen parameters, Tverberg's theorem generalizes to the topological setting, also known as Topological Tverberg's theorem.
%Namely, if $r$ is a power of a prime any continuous map from an $N$-dimensional simplex $\Delta_N$ to  $\mathbb{R}^d$, for $N\ge (d + 1)(r-1)$, maps $r$ pairwise disjoint faces of $\Delta$ so that they intersect in a common point.
If $r$ is a prime power, then the Tverberg theorem admits a topological generalization, known as the \emph{Topological Tverberg theorem}:
\begin{theorem} 
\label{thm:topTverberg}
%For every $d\in \mathbb{N}$, prime power $r$ and continuous map $\mu:\Delta_{(d+1)(r-1)}\to \R^d$, there exist $r$ pairwise disjoint faces $Y_1,\ldots, Y_r\subseteq \Delta$ such that $f(Y_1)\cap \ldots \cap  f(Y_r) \ne \emptyset$.
If $d\in \mathbb{N}$ and if $r$ is a prime power, then for every continuous map from the $(d+1)(r-1)$-dimensional simplex $\Delta_{(d+1)(r-1)}$ to $\R^d$, there exist $r$ pairwise disjoint faces of $\Delta_{(d+1)(r-1)}$ whose images intersect in a common point. 
\end{theorem}
This result was first proved in the case when $r$ is a prime by {B{\'a}r{\'a}ny, Shlosman and Sz{\"u}cs~\cite{BSS1981_topTverbergPrime} and later extended to prime powers by \"{O}zaydin~\cite{O1987_topTverberg}. On the other hand, Theorem~\ref{thm:topTverberg} is false if $r$ is not a prime power: By work of Mabillard and Wagner \cite{MabillardWagner2015} and a result of \"{O}zaydin~\cite{O1987_topTverberg}, a closely related result (the \emph{generalized Van Kampen--Flores theorem}) is false whenever $r$ is not a prime power and, as observed by Frick~\cite{F15_approx}, the failure of  Theorem~\ref{thm:topTverberg} for $r$ not a prime power follows from this by a reduction due to Gromov \cite{Gromov} and to Blagojevi\'{c}, Frick, and Ziegler \cite{BFZ_Constraints_2014}. The lowest dimension in which counterexamples are known to exist is $d=2r$~\cite{AMSW15_codim2,MabillardWagner2015}.
%We note that Theorem~\ref{thm:topTverberg} does not extend to all natural numbers $r$ due to a work of Wagner and Mabillard~\cite{MW14_elim} and \"{O}zaydin~\cite{O1987_topTverberg} as observed by Frick~\cite{F15_approx}.
%Nevertheless, such an extension could still be  possible for $d\le 12$, since the counterexamples are known to exist
%only in dimension 13 and higher~\cite{AMSW15_codim2},
We refer to the recent surveys~\cite{Barany2018,BZ2017_story,S18_user,Z18} for more background on the Topological Tverberg theorem and its history.

In the same vein, it is natural to wonder if our Theorem~\ref{thm:main} also extends to the topological setting. The straightforward approach to generalize our result would be to use the Topological Tverberg theorem instead of the Tverberg theorem and then keep fixing the pairs of simplices whose boundaries do not mutually intersect. To this end we need an adaptation of  Lemma~\ref{lem1}  and the fixing procedure
to the topological setting. The rest of our argument is free of any geometry except for the use of Carath\'{e}odory's theorem which is not really crucial. While extending  Lemma~\ref{lem1} is easy, showing the termination of the fixing procedure appears to be quite difficult except under the scenario which we discuss below.
First, we discuss planar extensions of Theorem~\ref{thm:main}.
An \emph{arrangement of pseudolines} $\mathbb{P}$ is a finite set of not self-intersecting open arcs, called \emph{pseudolines}, in $\mathbb{R}^2$ such that (i) For every pair $P_1,P_2\in \mathbb{P}$ of two distinct pseudolines, $P_1$ and $P_2$ intersect transversely in a single point, and (ii) $\mathbb{R}^2\setminus P$ is not connected for every $P\in \mathbb{P}$.
A drawing of a complete graph on $n$ vertices $K_n$ in the plane is \emph{pseudolinear} if the edges can be extended to an {arrangement of pseudolines}.


In the plane, it is not hard to see that  Theorem~\ref{thm:main} and its proof almost extends to the setting of  pseudolinear drawings of complete graphs. Since the Topological Tverberg theorem is only valid for prime powers $r$, the number of pairwise crossing triangles is slightly smaller than the number of vertices divided by 3.

\begin{theorem} \label{thm:mainTopPlane}
In a pseudolinear drawing of a complete graph  $K_{3n}$  we can find $m=(1-o(1))n$ vertex-disjoint and pairwise crossing triangles. Moreover, the topological discs bounded by these triangles intersect in a common point.
\end{theorem}

\begin{proof}
Let $\mathcal{D}$ be a pseudolinear drawing of  $K_{3n}$.
We put $m$ to be the largest prime power not larger than  $n$.
By the asymptotic law of distribution of prime numbers  $m=(1-o(1))n$.
Next, apply the Topological Tverberg theorem with $r=m$ and $d=2$ to a map $\mu:\Delta_{3m-3}\rightarrow \mathbb{R}^2$
which extends $\mathcal{D}$  as follows.
We define $\mu$ on the 1-dimensional skeleton of $\Delta_{3m-3}$ as  a restriction of $\mathcal{D}$  to some $K_{3m-2}$.
Note that every triangle of $K_{3m-2}$ is drawn by $\mathcal{D}$ as a closed arc without self intersections. The map $\mu$ extends to the 2-dimensional skeleton of $\Delta_{3m-3}$ so that every  2-dimensional face is mapped homeomorphically in $\mathbb{R}^2$. We  define the map $\mu$ on the rest of $\Delta_{3m-3}$ arbitrarily while maintaining continuity.

Analogously to the proof of Theorem~\ref{thm:main}, an application of the Topological Tverberg theorem gives us $m-1$ disjoint $2$-dimensional faces $F_1,\ldots, F_{m-1}$ of $\Delta_{3m-3}$ whose images under $\mu$ intersect in a common point.
 To this end we apply Carath\'eodory's theorem for drawings of complete graphs~\cite[Lemma 4.7]{BFK15_monotone} instead of the original version of Carath\'eodory's theorem.
Let $T_i$ denote the boundary of $F_i$, for $i=1,\ldots, m-1$.
Note that each $T_i$ is a triangle in $K_{3m-2}$.
If all pairs $T_i$ and $T_j$ are crossing then we are done. Otherwise, we perform the fixing operations, which can be done since  Lemma~\ref{lem1} easily extends to the setting in which  we replace simplices  by images of 2-dimensional faces of $\Delta_{3m-3}$  under $\mu$.

The procedure of applying successively the fixing operation terminates due to the following argument.
First, observe that 
every four vertices in a pseudolinear drawing of a complete graph induce
either a crossing free drawing of $K_4$ or a drawing of $K_4$ with exactly one pair of crossing edges in the interior of a disc bounded by a crossing free 4-cycle, see Figure~\ref{fig:pseudo} for an illustration.
It follows that if $\mu(T_i) \cap \mu(T_j)=\emptyset$, and let's say $\mu(F_i)\subset \mu(F_j)$, then the restriction of $\mu$ to the subgraph of $K_{3m-2}$  induced by $V(T_i)\cup V(T_j)$ is contained in $\mu(F_j)$.
Indeed, otherwise there exists a vertex $u\in V(T_i)$ and $v\in V(T_j)$ such that $\mu(uv)$ crosses an edge of $T_j$, let us denote it by $wz$, that is not incident to $v$. Then the restriction of $\mu$ to the subgraph of $K_{3m-2}$ induced by $\{u,v,w,z\}$ is a drawing of $K_4$ that is not pseudolinear  (contradiction). Therefore the volume argument goes through if we consider, say, volumes of $\mu(F_i)$'s.
\end{proof}


\begin{figure}
\centering
\includegraphics[scale=1]{pseudolinear}
\caption{The allowed drawings of $K_4$ in a pseudolinear drawing of $K_n$ (left). The forbidden drawing of $K_4$ in a pseudolinear drawing of $K_n$ (right).\label{fig:pseudo}}
\end{figure}

A drawing of a  graph in the plane is \emph{simple} if every pair of edges intersect at most once either at a common end point or in a proper crossing.
Clearly, all pseudolinear drawings of complete graphs are also simple, but not vice-versa, as illustrated in the last drawing in Figure~\ref{fig:pseudo}. Hence, it might be worthwhile to extend   Theorem~\ref{thm:mainTopPlane} to simple drawings of complete graphs.
 
If we want the interiors of the triangles to be pairwise intersecting, we only known that we can take $m$ at least $\Omega(\log^{1/6}n)$ which is easily derived from the following result of Pach, Solymosi and T\'oth~\cite{PST2003_unavoidable}. Every simple drawing of $K_n$ contains a drawing of $K_m$ that is weakly isomorphic to  a so-called convex  complete graph or a twisted complete graph, see Figure~\ref{fig:twisted}, for which Theorem~\ref{thm:mainTopPlane} holds. We omit the proof of the latter which is rather straightforward. For example, if in a twisted drawing of $K_{3n}$ the vertices are labeled as indicated in the figure, $\{\{0+i,n+i,2n+i\}| \ i=0,\ldots, n-1\}$ is a crossing Tverberg partition.

 If we do not insist on the interiors of the triangles to be pairwise intersecting, we know that $m$ can be taken to be at least $\Omega(n^{\varepsilon})$ for some small $\varepsilon>0$ by the following result of Fox and Pach~\cite{FP2012_coloring}.
Every simple drawing of $K_n$ contains $\Omega(n^{\varepsilon})$ pairwise crossing edges for some $\varepsilon>0$.

\begin{figure}

\centering
\includegraphics[scale=0.7]{twisted}
\caption{The twisted drawing of $K_6$. The bold grey and black edges form boundaries of two crossing triangles. The symbol x marks a  point contained in the interior of these two triangles.\label{fig:twisted}}
\end{figure}

Theorem~\ref{thm:mainTopPlane} could be generalized to hold for an appropriate high dimensional analog of pseudolinear drawings. Since this would require  introducing  many technical terms and would not offer substantially interesting content we refrain from doing so.

%In this section, we sketch the proof of the topological version of Theorem~\ref{thm:main}.
%Let us denote by $\Delta_n$ the simplex with the vertex set $[n]$. We call a continuous map $f: \Delta_n\to \R^d$ {\it simple} if, for any $X, Y\subset {[n]\choose d}$ such that $Y\cap X\ne \emptyset$, we have $f(\|Y\|)\cap f(\|X\|)=f(\|X\cap Y\|)$. (We use the name ``simple'' by analogy with simple topological drawings.)
%
%\begin{theorem} \label{thm:main}
%Let $f: \Delta_{(d+1)n-1}\to \R^d$ be a simple continuous map. Then we can find $m=(1+o(1))n$ vertex-disjoint $d$-faces $X_1,\ldots, X_m$ of $\Delta_{(d+1)n-1}$ such that the boundaries of $X_1,\ldots, X_m$ pairwise intersect.
%\end{theorem}
%
%We put $m$ to be the largest prime power that is strictly smaller than to $n$. Known results on the distribution of primes imply that $m=(1+o(1))n$. Next, apply the topological Tverberg theorem, stating that
%\begin{quote}
%  For all integer $d$, prime power $r$ and continuous map $\Delta_{(d+1)(r-1)}\to \R^d$, there are $r$ pairwise disjoint faces $Y_1,\ldots, Y_r\subset \Delta$ such that $f(Y_1)\cap \ldots f(Y_r)\ne \emptyset$.
%\end{quote}
%Due to the fact that $f$ is simple, the boundary of $f(X)$ for any $d$-face $X$ has a well-defined interior. Let $f':\Delta_{(d+1)n-1}\to \R^d$ be a continuous function that coincides with $f$ on all $\le (d-1)$-faces of $\Delta_{(d+1)n-1}$ and that maps the interior of $X$ onto the interior of the boundary of $f(X)$ for every $l$-face $X$, $l\ge d$.
%
%Choose any subset $Z\subset [n]$ of size $(d+1)(m-1)+1$ and apply the topological Tverberg theorem to the simplex formed on $X$ and the restriction of $f'$ on $X$.  Using the same argument as before, obtain images of $m$ disjoint $d$-faces that all contain the same point, w.l.o.g. $\nula$, in their interior. (Note that the ``interior'' here is with respect to the function $f$.)
%
%If all pairs of simplices are crossing then we are done. Otherwise, we perform the switching operations. We can do it since the analogue of Lemma~\ref{lem1} is valid for
%$$\mathcal C:=\big\{X\in{Z\choose d+1}:\ \text{the interior of the boundary of }f(\|X\|)\ \text{contains }\nula \big\}$$
%  (and the rest of the argument did not depend on any other information).
%
%
%Finally, {\it Fixing Pairs} terminate in this case as well since we go down in the lexicographical order on the volumes of the simplices, due to the simplicity of the map. Indeed, fix two disjoint $(d+1)$-element subsets $X_1$, $X_2$ of $Z$, such that $f'(\|X_1\|)$ contains $f'(\|X_2\|)$. Consider any other $(d+1)$-set $X_3\subset X_1\cup X_2$. Then $f'(\|X_3\|)\subsetneq f'(\|X_1\|)$ due to simplicity.

\subsection{Stronger conditions on crossings}\label{sec:stronger}
%One may ask if in Theorem~\ref{thm:main}, in addition to the condition that the boundaries of all $(d+1)$-sets of the partition intersect, it is possible to get a stronger property that $d$-boundaries of any two sets in the partition intersect. This condition automatically holds for  pairs of parts of size at most $d$, since they are equal to their $d$-boundary. However, it is not automatic for pairs $X_i,X_j$, where $X_i$ is of size $d+1$ and $X_i$ is of size $\le d$, since the latter does not necessarily intersect the $d$-boundary of the former. 
%
%It turns out that, in general, we cannot guarantee this property. Indeed, consider a set of $4$ points on the plane, which consists of a triangle and a point in its interior. Then the only Tverberg partition of this set has two parts: the outer triangle and the point in the interior. But the latter obviously does not intersect the boundary of the former.

A pair of vertex disjoint $(\lceil d/2 \rceil-1)$-dimensional simplices in general position in $d$-space does not intersect. 
Hence, the pairwise intersection of the boundaries of $\conv(X_i)$'s in the conclusion of   Theorem~\ref{thm:main}, cannot be strengthened to the pairwise intersection of lower than
$\lceil d/2 \rceil$-dimensional skeleta of the boundaries. 

Nevertheless, for $d=3$ one may ask if in the setting of Lemma~\ref{lem:unnesting}, we get a stronger property along the following lines. 
Can we guarantee the existence of a pair of vertex-disjoint tetrahedra $\{S,S'\}$ that both contain the origin and such that the boundary of a 2-dimensional face $F$  of $S$ is linked with the boundary of a 2-dimensional face $G$ of $S'$, meaning that the boundary of $F$ intersects $G$ and the boundary of $G$ intersects $F$? 
Again, the answer to this question is negative. Stefan Felsner and Manfred Scheucher (personal communication) found the following set of 8 points:
      $$(3, -2, 2),\ \ \ \
      (2, -5, 3),\ \ \ \
      (-3, 0, -4),\ \ \ \
      (-1, 2, 0),$$
      $$(1, -5, -4),\ \ \ \
      (4, 1, -2),\ \ \ \
      (-2, -5, -4),\ \ \ \
      (-3, 1, 3).$$
     This set determines a pair of disjoint tetrahedra both containing the origin $(0,0,0)$, but no two disjoint linked tetrahedra
     both containing the origin.

% \rado{we could perhaps draw this counterexample somehow 
% (Can anybody do this?) and state  which is the pair of disjoint tetrahedra both containing the origin $(0,0,0)$.  Do you see this fast? I am not sure if I can make anything of Stephan's suggestion (see his email from August 15) about viewing the point set in sage since I am not a sage user and his explanation of how the counterexample was found}
The example was found using a SAT solver who found an abstract order type with the required property. A realization of the order type with actual points was obtained with a randomized procedure.

\subsection{Computational complexity of finding a crossing Tverberg partition}
A natural question is whether we can find the partition of the point set given by Theorem~\ref{thm:main} efficiently, i.e.,  in polynomial time in the size of $X$. 
A straightforward way to construct an algorithm is to  make the proof of Theorem~\ref{thm:main} algorithmic.
To this end we first need an algorithm for finding a Tverberg partition and also a Tverberg point. Unfortunately, it is neither known whether this can be done in polynomial time, nor are there hardness results~\cite[Section 3.4]{TverbergComplexity18}. Since a Tverberg partition always exists, the decision problem is trivial, so NP-completeness cannot apply. It might still be possible to prove completeness of the problem for a suitable subclass of TFNP (containing search problems for which a solution is guaranteed to exist). It has been shown that the problem is contained in two such subclasses that have complete problems, namely PPAD and PLS~\cite[Theorem 4.9]{MMSS17_endofline}. While this makes it unlikely that the problem is complete for either of them, it could well be contained in and turn out to be complete for a subclass of PPAD$\cap$PLS, a number of which have been introduced and studied in recent years; see~\cite{DBLP:journals/corr/abs-1811-03841} and the references therein.

% Unfortunately, several results suggest that 
% an efficient algorithm for this problem is rather unlikely to exist. Since a Tverberg partition is guaranteed to exist, NP complexity theory does not apply to the algorithmic problem of finding it. Instead of NP, the PPAD complexity theory~\cite{P94_complexity} applies as the problem is known to be in PPAD~\cite[Theorem 4.9]{MMSS17_endofline}.

The only closely related hardness result we are aware of is the one by Teng~\cite[Theorem 8.14]{T92_points}, who proved  that checking whether a given point is a Tverberg point of a given point set is NP-complete.

A line of research on finding an approximate Tverberg partition efficiently was initiated by Miller and Sheehy~\cite{MS10_approx} and further developed in~\cite{MW13_approx,RS16_approx}.
In particular, Mulzer and Werner~\cite{MW13_approx} showed that it is possible to find  in time $d^{O(\log d)|X|}$ an approximate Tverberg partition of size $\left\lceil \frac{|X|}{4(d+1)^3}\right\rceil$, whereas Theorem~\ref{thm:Tverberg} guarantees the partition of size $\left\lceil \frac{|X|}{d+1}\right\rceil$.


If we aim only at an approximate algorithmic version of our Theorem~\ref{thm:main} along the lines of the result of Mulzer and Werner, we face the problem of efficiently fixing an (approximate) Tverberg partition to make it crossing. Due to the fact that our termination argument for the iterated \emph{Fixing Pairs} procedure relies on progress in the lexicographical ordering of the simplex volumes, we may potentially need exponentially many (in the size of~$X$) iterations before we arrive at a crossing partition. %We were unable to prove an upper bound better of than $2^n$ operations.
We leave it as an interesting open problem to prove or disprove that there is always a way to invoke the \emph{Fixing Pairs} operation only polynomially (or least subexponentially) many times in order to arrive at a partition required by Theorem~\ref{thm:main}. 


%%
%% Bibliography
%%

%% Please use bibtex, 

\bibliography{paper_references}


\end{document}
